\subsection{Basis-Resolved Admissibility}
  \label{subsec:basis-resolved-admissibility}

  \begin{proposition}[Basis-resolved admissibility on $Q_8$]
    \label{prop:basis-resolved}
    In the canonical Peter--Weyl basis of $\ell^2(Q_8)$, the maximal admissible
    effective amplitudes satisfy
    \begin{equation}
      A_{\rho_4}^{\max}
      =
      \frac{c_\chi}{\sqrt{6}}
      >
      \frac{c_\chi}{\sqrt{8}}
      =
      A_{\rho_k}^{\max},
      \qquad
      k=1,2,3,
    \end{equation}
    whereas the corresponding maximal raw amplitudes satisfy
    \begin{equation}
      a_{\rho_4}^{\max}
      =
      \frac{c_\chi}{\sqrt{6}\,\|\psi_{\rho_4}\|_\infty}
      =
      \frac{2}{\sqrt{6}}\,c_\chi
      \approx 0.816\,c_\chi,
    \end{equation}
    \begin{equation}
      a_{\rho_k}^{\max}
      =
      \frac{c_\chi}{\sqrt{8}\,\|\psi_{\rho_k}\|_\infty}
      =
      c_\chi,
      \qquad
      k=1,2,3.
    \end{equation}
    Thus the non-abelian sector is favoured in effective amplitude, but not in raw
    amplitude.
  \end{proposition}

  \begin{proof}
    The effective-amplitude statement follows from
    $A_n^{\max} = c_\chi/\sqrt{\lambda_n}$ together with
    $\lambda_{\rho_4} = 6$ and $\lambda_{\rho_k} = 8$.
    The raw-amplitude statement follows from
    $a_n^{\max} = A_n^{\max}/\|\psi_n\|_\infty$, using
    $\|\psi_{\rho_4}\|_\infty = 1/2$ and
    $\|\psi_{\rho_k}\|_\infty = 1/\sqrt{8}$.
  \end{proof}
